%% ====================================================================
%% Command line basics
%% ====================================================================

\begin{frame}{Why learn the command line?}
Your sequencing run is back:
\begin{itemize}
    \item 200 FASTQ files, 500 GB in total
    \item A laptop with 8 GB of memory and a nearly full disk
    \item An assembler with no buttons: only a command line
\end{itemize}
\bigskip
\textbf{Today you will learn to}
\begin{itemize}
    \item Find, check and process files from the command line
    \item Keep your scripts safe with Git
    \item Install tools yourself with Conda
    \item Run your analysis on the HPC with Slurm
\end{itemize}
\end{frame}

\section[CLI]{Command-line basics}
\begin{frame}{Before we start: open a terminal}
\begin{cmdtable}{ll}
    \textbf{Computer} & \textbf{How to open a terminal}                \\ \midrule
    Mac               & \cmd{Cmd+Space}, type \cmd{Terminal}, \cmd{Enter} \\
    Windows           & Start menu: \textbf{Ubuntu} (WSL)              \\
    Linux             & \cmd{Ctrl+Alt+T}                               \\
    HPC               & A terminal in Open OnDemand (browser)          \\
\end{cmdtable}
\bigskip
You type commands in the terminal and press \cmd{Enter} to run them

\medskip
\hint{Open OnDemand: \link{https://hpcs05.op.umcutrecht.nl}.
    Windows: install WSL before the workshop}
\end{frame}

\begin{frame}[fragile]{Before we start: get the data for today}
Git clone the workshop repository
\begin{lstlisting}
cd ~                          # laptop: home
cd /hpc/<group>/<user>/data   # HPC: data directory
git clone https://github.com/bayraktar1/HPC_intro.git
cd HPC_intro/exercise_data
\end{lstlisting}
\begin{itemize}
    \item On the HPC, keep data and environments out of your home directory (7 GB)
    \item \cmd{git} not found? Raise your hand
\end{itemize}
\end{frame}

\begin{frame}{How the exercises work}
\begin{itemize}
    \item The exercises are in \cmd{HPC_intro/exercises}, one file per section.
          Easiest to read on GitHub: \href{https://github.com/bayraktar1/HPC_intro}{\cmd{github.com/bayraktar1/HPC_intro}}
    \item When a slide says \textbf{Exercises}, do the numbers it lists
    \item Every exercise has a level:
\end{itemize}
\medskip
\hspace*{\leftmargini}\begin{tabular}{@{}l@{\hspace{2em}}l@{}}
    \level{1} & everyone does these           \\
    \level{2} & do these if you can           \\
    \level{3} & optional, if you finish early \\
\end{tabular}
\end{frame}

\subsection{Files and paths}
\begin{frame}{Every command has the same anatomy}
\diagram{command-anatomy}

\bigskip
\begin{itemize}
    \item \textbf{Options} start with \cmd{-} or \cmd{--} and change what a command does
    \item \textbf{Arguments} tell it what to work on: a file, a directory
    \item \cmd{<...>} on these slides: fill in your own value
    \item Stuck?
          \cmd{man ls} (press \cmd{q} to quit) or \cmd{ls --help}
\end{itemize}
\end{frame}

\begin{frame}{Moving around: where am I and what is here?}
\small
\begin{cmdtable}{ll}
    \textbf{Command}       & \textbf{What it does}                                 \\
    \midrule
    \cmd{pwd}              & Print where you are (working directory)               \\
    \cmd{ls}               & List what is here                                     \\
    \cmd{ls -lh <dir>}     & List with details and readable sizes                  \\
    \cmd{ls -a}            & Also show hidden files: names that start with \cmd{.} \\
    \cmd{cd <dir>}         & Go into a directory                                   \\
    \cmd{cd ..}            & Go one level up                                       \\
    \cmd{cd ~} or \cmd{cd} & Go to your home directory                             \\
    \cmd{mkdir <dir>}      & Make a directory                                      \\
\end{cmdtable}
\end{frame}

\begin{frame}[fragile]{Paths: how to get from here to there}
\begin{columns}[T]
    \column{0.38\textwidth}
    \begin{lstlisting}[style=output]
~  (your home)
`-- HPC_intro
    `-- exercise_data <-
        |-- sample_A
        |-- sample_B
        `-- sample_C
            `-- old
\end{lstlisting}
    \column{0.6\textwidth}
    \begin{lstlisting}[style=footnotesize]
# start in exercise_data
cd sample_C       # down into sample_C
cd old            # further down
cd ../..          # two levels up
cd sample_A       # down into sample_A
cd ../sample_B    # up, into sample_B
cd ~              # home (or just: cd)
cd ~/HPC_intro/exercise_data
                  # works from anywhere
\end{lstlisting}
\end{columns}
\medskip
\small
\begin{itemize}
    \item \cmd{..} is one level up, \cmd{.} is where you are now
    \item A path that starts with \cmd{/} or \cmd{~} works from anywhere.
          Any other path starts from where you are now: check with \cmd{pwd}
\end{itemize}
\hint{On the HPC, \cmd{HPC_intro} is in \cmd{/hpc/<group>/<user>/data} instead of \cmd{~}}
\end{frame}

\begin{frame}{Exercises: moving around}
\begin{exlist}{1_command_line.md}
    \ex{1.1}{Where am I?}{1}
    \ex{1.2}{List and look}{1}
    \ex{1.3}{Down and back up}{2}
    \ex{1.4}{Your own directory}{1}
\end{exlist}
\end{frame}

\begin{frame}{Keyboard tricks that save you time}
\small
\begin{cmdtable}{ll}
    \textbf{Keys}             & \textbf{What they do}           \\
    \midrule
    \cmd{Tab}                 & Complete a file or command name \\
    \cmd{Tab Tab}             & Show all possible completions   \\
    $\uparrow$ / $\downarrow$ & Previous / next command         \\
    \cmd{Ctrl+R}              & Search your previous commands   \\
    \cmd{Ctrl+C}              & Stop the running command        \\
    \cmd{q}                   & Quit \cmd{less} and \cmd{man}   \\
    \cmd{Ctrl+L}              & Clear the screen                \\
\end{cmdtable}

\bigskip
Use \cmd{Tab} all the time: it is faster and prevents typos in file names
\end{frame}

\begin{frame}{Looking at files}
\small
\begin{cmdtable}{ll}
    \textbf{Command}         & \textbf{What it does}                 \\
    \midrule
    \cmd{cat <file>}         & Print a whole file                    \\
    \cmd{head -n <N> <file>} & Print the first \cmd{N} lines         \\
    \cmd{less <file>}        & Scroll through a file (\cmd{q} quits) \\
    \cmd{wc -l <file>}       & Count the lines                       \\
\end{cmdtable}

\bigskip
These only read: they never change your file
\end{frame}

\begin{frame}{Exercises: looking at files}
\begin{exlist}{1_command_line.md}
    \ex{1.5}{Look inside a file}{1}
    \ex{1.6}{Count lines}{2}
\end{exlist}
\end{frame}

\begin{frame}{Changing files}
\small
\begin{cmdtable}{ll}
    \textbf{Command}       & \textbf{What it does}                                             \\
    \midrule
    \cmd{nano <file>}      & Edit a file (\cmd{Ctrl+O}, \cmd{Enter}: save; \cmd{Ctrl+X}: quit) \\
    \cmd{cp <file> <copy>} & Copy a file                                                       \\
    \cmd{mv <file> <new>}  & Move or rename a file                                             \\
    \cmd{rm <file>}        & Delete a file: \alert{no undo, no trash!}                         \\
    \cmd{rmdir <dir>}      & Remove an empty directory                                         \\
\end{cmdtable}
\end{frame}

\begin{frame}{Exercises: changing files}
\begin{exlist}{1_command_line.md}
    \ex{1.7}{Write notes}{1}
    \ex{1.8}{Copy, rename, delete}{2}
\end{exlist}
\end{frame}

\begin{frame}{Some characters have a special meaning}
\small
\begin{cmdtablex}{lLl}
    \textbf{Char} & \textbf{Meaning}                                     & \textbf{Example}               \\
    \midrule
    \cmd{~}       & Your home directory                                  & \cmd{cd ~}                     \\
    \cmd{.}       & The current directory                                & \cmd{cp ~/notes.txt .}         \\
    \cmd{..}      & One directory up                                     & \cmd{cd ..}                    \\
    \cmd{/}       & Between directories, or the top                     & \cmd{/home/student}            \\
    \cmd{*}       & Any text (wildcard)                                  & \cmd{ls sample_*}              \\
    \cmd{"..."}   & Keep text with spaces together                       & \cmd{mkdir "my results"}       \\
\end{cmdtablex}
\end{frame}

\begin{frame}[fragile]{Checksums tell you if your data arrived intact}
\small
\begin{itemize}
    \item A checksum is a file's fingerprint: change one byte and it changes completely
    \item Transfers can silently corrupt files: a broken \cmd{.fastq.gz} can look fine until your analysis fails
    \item Data often comes with checksums: check them \textbf{before} you start, and before deleting originals
\end{itemize}
\begin{lstlisting}
md5sum sample_A/sample_A_R1.fastq.gz   # one file
md5sum -c MD5.txt                      # check the list
\end{lstlisting}
\begin{lstlisting}[style=output]
sample_A/sample_A_R1.fastq.gz: OK
sample_A/sample_A_R2.fastq.gz: OK
...
md5sum: WARNING: 1 computed checksum did NOT match
\end{lstlisting}
\hint{A file marked \cmd{FAILED}: transfer it again}
\end{frame}

\begin{frame}{Exercises: checksums}
\begin{exlist}{1_command_line.md}
    \ex{1.9}{Check the reads}{1}
    \ex{1.10}{Compare by hand}{2}
\end{exlist}
\end{frame}

\subsection{Searching and automating}
\begin{frame}[fragile]{More powerful commands: \cmd{find}}
\begin{block}{Syntax}
    \cmd{find <where> -name "<pattern>"}
\end{block}
\begin{lstlisting}
find . -name "*.fasta"                 # here and below
find sample_C -type f -name "*.fasta"  # only files
\end{lstlisting}
\hint{Quote the pattern, so the \cmd{*} goes to \cmd{find} instead of being filled in by the shell}
\end{frame}

\begin{frame}[fragile]{More powerful commands: \cmd{grep}}
\begin{block}{Syntax}
    \cmd{grep [options] "<pattern>" <file>}
\end{block}
\begin{lstlisting}
grep "^>" sample_B/assembly.fasta       # header lines
grep -c "^>" sample_B/assembly.fasta    # count them
grep -i "chromosome" sample_B/assembly.fasta
\end{lstlisting}
\small
\begin{itemize}
    \item \cmd{^} in a pattern means: the start of the line.
          \cmd{-i} ignores upper/lower case
    \item \alert{Always quote the \cmd{>}}: \mbox{\cmd{grep -c > assembly.fasta}} empties the file!
\end{itemize}
\end{frame}

\begin{frame}[fragile]{Pipes send the output of one command to the next}
\cmd{<command> | <command>}: everything the left command prints goes into the right one
\begin{lstlisting}
ls sample_B | wc -l          # count files

# the first 3 headers
grep "^>" sample_C/assembly.fasta | head -n 3

# > writes the output to a file instead of the screen
grep "^>" sample_C/assembly.fasta > headers.txt
\end{lstlisting}
{\small You can chain as many commands as you like: \cmd{a | b | c}\par}

\medskip
\begin{alertblock}{\cmd{>} overwrites without asking}
    \cmd{head -n 3 a.fasta > a.fasta} leaves you with an \textbf{empty} \cmd{a.fasta}:
    the shell empties the file before \cmd{head} even runs.
    Always write to a new file.
\end{alertblock}
\end{frame}

\begin{frame}{Exercises: find and grep}
\begin{exlist}{1_command_line.md}
    \ex{1.11}{Find files}{1}
    \ex{1.12}{Count sequences}{2}
    \ex{1.13}{Count files with a pipe}{2}
    \ex{1.14}{Plasmid headers}{3}
\end{exlist}
\end{frame}

\begin{frame}[fragile]{More powerful commands: \cmd{for} loop}
\begin{block}{Syntax}
    \cmd{for <var> in <list>; do <command> "$<var>"; done}
\end{block}
\begin{lstlisting}
for f in sample_B/*.fastq.gz; do
    md5sum "$f"
done
\end{lstlisting}
\pause
{\small This runs the command once for every file:\par}
\begin{lstlisting}
# round 1: $f is sample_B/sample_B_R1.fastq.gz
md5sum "sample_B/sample_B_R1.fastq.gz"
# round 2: $f is sample_B/sample_B_R2.fastq.gz
md5sum "sample_B/sample_B_R2.fastq.gz"
\end{lstlisting}
\pause
\begin{itemize}
    \item \cmd{f} is a variable: every round it holds the next file name
    \item \cmd{$f} gives its value.
          Quote it, \cmd{"$f"}, so names with spaces do not break
\end{itemize}
\end{frame}

\begin{frame}[fragile]{Test a loop with \cmd{echo} before you run it}
Put \cmd{echo} in front: the loop prints the commands instead of running them
\begin{lstlisting}
for f in sample_*/*.fastq.gz; do echo md5sum "$f"; done
\end{lstlisting}
\pause
\begin{lstlisting}[style=output]
md5sum sample_A/sample_A_R1.fastq.gz
md5sum sample_A/sample_A_R2.fastq.gz
md5sum sample_B/sample_B_R1.fastq.gz
md5sum sample_B/sample_B_R2.fastq.gz
md5sum sample_C/sample_C_R1.fastq.gz
md5sum sample_C/sample_C_R2.fastq.gz
\end{lstlisting}
\pause
Looks right?
Remove the \cmd{echo} to run it for real:
\begin{lstlisting}
for f in sample_*/*.fastq.gz; do md5sum "$f"; done
\end{lstlisting}
\end{frame}

\begin{frame}[fragile]{Scripts: save your commands in a file}
A script is commands saved in a file, so you can run them again
\begin{lstlisting}
#!/bin/bash
# check.sh: check all reads against MD5.txt
md5sum -c MD5.txt
\end{lstlisting}
\begin{lstlisting}
chmod +x check.sh   # make it executable (once)
./check.sh          # run it: ./ means "the file here"
\end{lstlisting}
\begin{itemize}
    \item The first line, \texttt{\#!/bin/bash}, says: run this file with bash
\end{itemize}
\end{frame}

\begin{frame}{Exercises: scripts and loops}
\begin{exlist}{1_command_line.md}
    \ex{1.15}{Your first script}{2}
    \ex{1.16}{A for loop}{3}
    \ex{1.17}{A loop in a script}{3}
\end{exlist}
\end{frame}

\subsection{Remote access}
\begin{frame}[fragile]{SSH lets you work securely on another computer}
\begin{itemize}
    \item SSH gives you a terminal on another computer, over an encrypted connection
    \item We use it to log in to the HPC, and Git uses it to talk to GitHub
\end{itemize}
\begin{lstlisting}
ssh <user>@<server>     # log in to a remote computer
exit                    # log out again
\end{lstlisting}
\end{frame}

\begin{frame}{SSH keys: a private key and a public key}
\diagram{ssh-keys}

\bigskip
\small
\begin{itemize}
    \item \cmd{ssh-keygen} creates a key pair: two files in \cmd{~/.ssh}
    \item The \textbf{public} key is like a lock: give it to the servers you want to use
    \item The \textbf{private} key opens it: keep it on your computer, protect it with a passphrase, \textbf{never share it}
    \item No more passwords: the server checks you have the matching private key
    \item On the HPC you can log in through Open OnDemand, without a key
\end{itemize}
\end{frame}
